\section{Related Works} \label{sec:related_works}

\subsection{Automated Stone Stacking}
Furrer et al.~\cite{furrer2017stacking} combined gradient-based pose search with physics simulation for stone stacking.
Johns et al.~\cite{johns2020autonomous} integrated perception, planning, and excavator manipulation for wall construction, later using constrained registration and signed-distance-field classification for prescribed wall geometries~\cite{johns2023framework}.
Liu et al.~\cite{liu2021drystacking} selected simulation-generated poses using hierarchical filters informed by masonry heuristics.
Liu and Napp~\cite{liu2023stability} incorporated structural strength evaluation into sequence planning with physical validation, later introducing stackability measures for corbelled shelter construction~\cite{liu2025clochain}.
Menezes et al.~\cite{menezes2021rocks} used deep Q-learning on elevation maps to select planar placements, rewarding target occupancy and penalizing object displacement.
Their fixed-sequence, fixed-orientation construction strategies inform our initial pose sampler.

\subsection{Irregular Object Packing}
Wang and Hauser~\cite{wang2022binpacking} packed nonconvex objects through heightmap minimization, considering stability, robot motion feasibility, and replanning for execution errors.
Huang et al.~\cite{huang2023binpacking} jointly planned object order and placement through hierarchical reinforcement learning with multi-view heightmaps and compactness, pyramidality, and simulated-stability rewards.

\subsection{Contact-Aware Pose Optimization}
Lee et al.~\cite{lee2025contactfactorgraph} use differentiable Contact Factor Graphs for gradient-based manipulation planning.
BODex~\cite{chen2025bodex} couples lower-level contact-force optimization with upper-level grasp synthesis, while GraspQP~\cite{zurbrugg2025graspqp} defines a differentiable force-closure energy through a quadratic program.
Our optimizer couples candidate placement with structure-wide equilibrium, using bodywise force decomposition for parallel evaluation.

\subsection{Simulation-Based Sequential Manipulation Planning}
King et al.~\cite{king2017unobservable} extended MCTS to nonprehensile rearrangement under uncertainty using learned and heuristic rollout policies.
Huang et al.~\cite{huang2022pmbs} introduced Parallel Monte Carlo Tree Search with Batched Simulations (PMBS), coupling parallel search with GPU rigid-body simulation for object retrieval from clutter.
Our stacking planner adopts batched evaluation, combining geometric sampling, stability-aware pose refinement, and structural simulation at each expansion.
